\section{First-sign Self-RAKL evidence}
The v2 paper expansion itself is another controlled self-application: literature routes were registered before assimilation, every added citation was mapped to a specific inherited mechanism or novelty correction, and the paper audit exposed that the end-to-end information lifecycle had been distributed across modules without one executable atomic trace. Round~43 therefore froze and implemented a 17-stage lifecycle plus a machine-receipted mini research trace before using those results in this manuscript. This is again same-context engineering evidence, not independent self-evolution.

Applying the method to itself has already exposed several real defects: same-context expert roles had initially been too easy to confuse with independent review; storage and prompt growth had not been sufficiently separated; metacognitive diagnosis did not yet determine how the system should learn; missing method operators had not been clearly separated from missing evidence or implementation defects; historical meta-fiber identifiers collided; and recent cross-domain model-discovery review exposed implicit model-criticism and assumption-sensitivity mechanics that needed explicit executable contracts. The canonical historical-ledger compiler also rejected incorrect Round-42 fiber allocations during this preprint preparation.

These are useful first-sign self-correction results because the framework's own controls caught and repaired problems in its representation and method. They are not strong self-evolution evidence: the work shares one overall orchestration context and has not yet executed the frozen fresh-assurance coding-agent bootstrap benchmark. When the current same-context evidence is evaluated against the strong bootstrap contract, missing registered development/fresh-assurance quantities correctly prevent the strong verdict rather than being silently inferred.


The Obsidian incident is another preserved first-sign failure. A relevant adjacent knowledge-navigation system was available in the world but not surfaced by the framework's earlier search routes; a later external prompt exposed the miss. Rather than adding the product name to a permanent exception list, \RAKL{} records the miss as failure memory and expands the route contract prospectively. Because the repair was designed after seeing the missed case, rediscovering Obsidian carries no fresh-transfer credit.

\subsection{What the existing Self-RAKL history can and cannot establish}

The same-context development history is valuable because it contains real mistakes, not only successful demonstrations. Framework iterations have exposed missing uncertainty propagation, evidence-lineage dependence, retrieval closure, authority leakage, measurement-basis ambiguity and the exogenous Obsidian miss. Repairs to those defects are evidence of \emph{self-correction under challenge}. They are not independent evidence that the method can autonomously discover arbitrary future weaknesses.

The distinction depends on chronology. Once a failure is visible, a repair tailored to it is development evidence. Replaying the repaired system on the same failure can test regression but cannot supply transfer credit. The next evidential step is a hidden or fresh realization that targets the same abstract capability without revealing the original fix. This is analogous to the difference between explaining a training example and generalizing a learned procedure.

The Obsidian incident is especially informative. The relevant neighboring concept was available in the external world but not discovered by the incumbent search routes until introduced exogenously. The correct response was not to add ``search for Obsidian'' to a special-case list. It was to identify the more general failure class, ontology-conditioned closure, and expand the prospective route contract. The repair will earn stronger credit only if it improves discovery on concepts not used to design it.

Accordingly, the present paper uses a hierarchy of self-evidence: defect discovery, causal attribution, implementation repair, regression validation, hidden development validation, and fresh assurance. Only the later levels support a transferable self-evolution claim, and even then only for the registered method surface and task family. This hierarchy prevents an attractive narrative of recursive improvement from outrunning the evidence.
